\subsection*{Finite trial positivity and the exact omitted-tail correction}

The trial-head term in Proposition~\ref{prop:v120-structured-uniform}
must be distinguished from the exact Schur complement. A finite-support
certificate can settle the former without settling the latter.

\begin{proposition}[Finite-support trial forms and Galerkin Schur bounds]
\label{prop:v121-finite-trial}
Let the outer head be finite dimensional, let $T\succeq cI$ with
$c>0$ be a positive self-adjoint tail operator, and let $B$ be the
bounded head-to-tail coupling. For $X$ with range in $\mathcal D(T)$,
define $K_X,r$ as in Proposition~\ref{prop:v120-structured-uniform}
and put $S_\infty=F-B^*T^{-1}B$. Then
\begin{align}
 K_X-S_\infty
  &=(X-T^{-1}B)^*T(X-T^{-1}B)\succeq0,
       \label{eq:v121-square}\\
 S_\infty&=K_X-r^*T^{-1}r.\label{eq:v121-residual}
\end{align}
Let $P_M$ be the orthogonal projection onto a finite tail space in $\mathcal D(T)$, and
write $T_M=P_MT|_{\operatorname{Ran}P_M}$, $B_M=P_MB$ and
$S_M=F-B_M^*T_M^{-1}B_M$. If $X=P_MX$, then
\begin{equation}\label{eq:v121-finite-trial}
 K_X=\begin{bmatrix}I\\-X\end{bmatrix}^{*}
       \begin{bmatrix}F&B_M^*\\B_M&T_M\end{bmatrix}
       \begin{bmatrix}I\\-X\end{bmatrix}.
\end{equation}
In particular, $S_M\succeq0$ implies $K_X\succeq0$ for every such
$X$, and strict positivity transfers as well. This is the complete
operator's $K_X$, not a truncation of that trial form.

For the exact finite solve $X_M=T_M^{-1}B_M$, embedded in the complete
tail, one has $K_{X_M}=S_M$ and
\begin{equation}\label{eq:v121-galerkin-defect}
 \boxed{S_\infty=S_M-r_M^*T^{-1}r_M,
 \qquad r_M=B-TX_M.}
\end{equation}
If the finite tail spaces are nested and their union is a form core
for $T$, then $S_M\downarrow S_\infty$ in matrix norm.
\end{proposition}
\begin{proof}
Expanding the first square gives $K_X-F+B^*T^{-1}B$.
Since $r=-T(X-T^{-1}B)$, its value is $r^*T^{-1}r$,
which proves both identities. In \eqref{eq:v121-finite-trial}
every trial lies in the head plus the retained finite tail space,
so compression changes none of the quadratic terms. Finite square completion gives
\[
 K_X=S_M+(X-T_M^{-1}B_M)^*T_M(X-T_M^{-1}B_M).
\]
Thus $K_X\succeq S_M$, and the graph map $h\mapsto(h,-Xh)$ is
injective. Substitution of the exact finite solve proves
\eqref{eq:v121-galerkin-defect}; its retained residual rows vanish,
but its omitted rows need not vanish.

For each head vector $h$, the variational identity is
\[
 \ip{B^*T^{-1}Bh}{h}
 =\sup_{y\in\mathcal D(T^{1/2})}
 \{2\Re\ip{Bh}{y}-\norm{T^{1/2}y}^2\}.
\]
Restricting the supremum to $\operatorname{Ran}P_M$ gives
$\ip{B_M^*T_M^{-1}B_Mh}{h}$. Nestedness and form-core density make
these quantities increase to the full supremum. Polarization and
finite dimensionality of the head give matrix-norm convergence of
the Schur complements in the stated decreasing order.
\end{proof}

A finite list of positive $S_M$ therefore does not prove positivity
of $S_\infty$. Positivity for every member of a cofinal sequence at
a fixed window would imply $S_\infty\succeq0$, but strict positivity
would still require a positive limiting margin. The structured inverse
bounds concern this same $T$; a comparison metric or a shifted lower
operator does not redefine the target Schur complement.

\begin{remark}[The semidefinite range condition]\label{rem:v121-range}
For $T\succeq0$, the notation $T^{-1}B$ requires an additional
condition. A form version is available if
$\operatorname{Ran}B\subset\operatorname{Ran}T^{1/2}$. Define the
reduced inverse $C=T^{-1/2}B$, with values orthogonal to $\ker T$.
For $X$ with range in the form domain, use
$(T^{1/2}X)^*T^{1/2}X$ in the definition of $K_X$. Then
\[
 K_X-(F-C^*C)=(T^{1/2}X-C)^*(T^{1/2}X-C)\succeq0.
\]
Orthogonality to $\ker T$, or membership in the closure of its range,
does not suffice when the range is not closed. At the certified
$\lambda=4$ split there is no such issue:
Proposition~\ref{prop:v120-lambda4-tail} gives
$T\succeq10^{-8}D\succeq1.7940\cdot10^{-8}I$ on $Q_{16}$.
\end{remark}

\begin{proposition}[A certified local zero head-error term]
\label{prop:v121-alpha-zero}
For the actual $\lambda=4$ Weil form with outer head $E_{16}$,
$K_X\succ0$ in both parity sectors for $X=0$ and for the frozen
dyadic solves supported on modes $17,\ldots,256$ specified in
Appendix~\ref{app:v121-audit}. Consequently $\alpha_4=0$ is
admissible in Proposition~\ref{prop:v120-structured-uniform} for
each of these choices, on the complete operator.
\end{proposition}
\begin{proof}
Assemble the exact finite entries with the separate logarithmic
diagonal, normalized parity basis and $L=2\log4$. Outward interval
$LDL^*$ factorization of $K_X$ verifies all $17$ even pivots and
all $16$ odd pivots strictly positive, both for $X=0$ and for the
specified dyadics. The identical witnesses pass at $768$ and $896$
bits; Appendix~\ref{app:v121-audit} gives the files and gates.
Real symmetric positivity also proves positivity on the complex
parity spaces. Equation~\eqref{eq:v121-finite-trial} identifies these
matrices with the complete trial forms. No infinite Schur sign or
small residual is used in this argument.
\end{proof}

This result does not certify $\alpha_\lambda=0$ at untested windows
or for larger unsupported solves. Positivity of an exact Schur
complement would guarantee $K_X\succeq0$ for every trial; without
that premise the least admissible $\alpha$ can depend on $X$.
For example, $F=0$, $B=T=1$ gives $S=-1$, $K_0=0$ and $K_1=-1$.
Even $K_0\succ0$ need not help the signed target: its residual is
the entire coupling $B$.

With a proved $\alpha_\lambda=0$ along a growing family, the precise
remaining sufficient estimates would be
\begin{equation}\label{eq:v121-scaled-target}
 e_\lambda\longrightarrow0,
 \qquad t_\lambda/\sqrt{\mu_\lambda}\longrightarrow0.
\end{equation}
The second condition is stronger than $t_\lambda\to0$ when
$\mu_\lambda\to0$. Geometric convergence of an inverse approximation
improves an enclosure of $r_\lambda^*T_\lambda^{-1}r_\lambda$;
it does not force that residual energy to vanish for a fixed trial.
These growing-window estimates remain unproved. Failure of this
sufficient bound would not itself refute Weil positivity.

